\section{Conclusion}\label{sec:8}

Linguistic transparency is best understood as a family of mediated accessibility relations rather than as a single syntactic property. The umbrella schema, \mention{$X$ is transparent with respect to property $P$ for relation $R$ when $P$ remains accessible through $X$ for the purposes of $R$}, is a comparative concept: an analyst's question, applicable to any language, that asks which configurations let an internal property remain available to a grammatical or interpretive relation. The schema doesn't claim that any language has a descriptive category called \term{transparency}. It states the question.

The conclusion developed in §\ref{sec:7} is that the umbrella schema locates \term{projectibility profiles} rather than naming a single kind. The $X$/$P$/$R$ values across the family are heterogeneous: no shared stabilising source unifies them, and the projectible content lives at the level of the individual profile, not the level of the schema. The schema's job is to make the projectibility question askable; profiles earn scientific use when they identify an observable base licensing non-trivial projection to a target, stabilised by an identifiable source.

The profiles the paper has worked through:

\begin{itemize}
\tightlist
\item
  \term{Number-transparent quantificational nouns} (§\ref{sec:2-4}). A small, relatively stable category in the descriptive system of CGEL English (qua linguistic class, qua maintained-and-projectible category in §\ref{sec:7-1}'s stronger sense, once maintenance is attributed). Class membership predicts agreement override, characteristic determiner patterns, partitive-vs-non-partitive \mention{of}-complement availability, behaviour under fused-head construction, count/non-count percolation, and premodification tendencies. Stabilised by conventional lexical inheritance. The strongest candidate to pass Slater's \autocite*{slater-2015-natural-kindness} broad stability-plus-projection diagnostic among the cases examined; the within-domain audit is open.
\item
  \mention{The form-meaning accessibility family} (§§\ref{sec:3}--\ref{sec:5}). Three within-domain projectibility profiles related by analytical theme rather than by demonstrated cross-domain projection: a construction-level profile stabilised by constructional inheritance (CxG); a compound-level profile stabilised by usage-based entrenchment (Bell and Schäfer); a derivationally-complex-word-level profile stabilised by processing architecture (Heyer and Kornishova). Each licenses within-domain projection. Whether the three compose into a cross-domain projectibility cluster on shared items, and whether Khalidi \autocite*{khalidi-2018-causal-networks}-style causal-network refinement applies to that cluster, depend on within-items cross-modal evidence the present paper doesn't provide.
\end{itemize}

Coordination is excluded as a non-instance (§\ref{sec:2-3}): nothing was going to be encapsulated. Referential transparency is acknowledged as the boundary case (§\ref{sec:4-3}) where the schema's mediator is a context rather than a form; the philosophical-semantics tradition is named without being developed.

What is left open. Korean case-stacking and scrambling (§\ref{sec:2-6}) are deferred for separate co-authored development; the comparative-concept frame licenses asking the schema's question of those patterns with $R$-values different from agreement (argument-role recovery, scope, information-structure assignment), but the language-particular content is pending. The within-items cross-modal evidence that would test whether the form-meaning family has Khalidi-style causal-hierarchical structure (corpus predictors, ratings, priming at multiple SOAs, and interpretation accuracy on a single item set) is future work. Acquisition trajectory is a plausible further extension consistent with usage-based theories but not directly grounded by the present cases.

The paper's closing claim. \term{Transparency} in linguistics names a family of mediated accessibility relations that the umbrella schema locates as candidate projectibility profiles. A transparency attribution is scientifically useful when it identifies a profile: an observable base that licenses non-trivial projection to a target, stabilised by an identifiable source. Some profiles support strong projection; others weak; some none. Slater's SPC framework and Khalidi's causal-network account are diagnostic tools for assessing how well a profile supports induction. Kindhood is the downstream verdict on profiles that earn it, not the master question of the analysis.
